\chapter{Kiến trúc hệ thống}
\label{chap:architecture}

\section{Tổng quan thành phần}

Hình~\ref{fig:architecture} tách hệ thống thành hai luồng có điểm nối. Luồng nghiệp vụ
Nova Bank phục vụ người dùng tương tác qua trình duyệt; Next.js chỉ trình bày
và gọi API, còn FastAPI xác thực, kiểm tra quyền và ghi dữ liệu vào SQLite.
Luồng giám sát nhận file log, chuẩn hóa chúng thành Event, thực thi các rule,
chấm điểm và ghi artifact. Cùng một repository RBAC được dùng làm nguồn sự
thật của quyết định quyền trong cả hai luồng. Audit adapter đọc bản ghi Nova
Bank theo thứ tự thời gian và xuất CSV đúng hợp đồng parser, tạo đường đi
end-to-end từ quyết định nghiệp vụ tới detector.

\begin{figure}[htbp]
\centering
\resizebox{\textwidth}{!}{\begin{tikzpicture}[
  node distance=8mm and 12mm,
  box/.style={draw, rounded corners, align=center, minimum height=10mm, minimum width=28mm, fill=blue!6},
  store/.style={draw, cylinder, shape border rotate=90, aspect=0.28, align=center, minimum height=13mm, minimum width=18mm, fill=orange!14},
  flow/.style={-Latex, thick}
]
\node[box] (web) {Next.js\\Nova Bank UI};
\node[box, right=of web] (api) {FastAPI\\API + RBAC enforcement};
\node[store, right=of api] (db) {SQLite\\RBAC, session, audit};
\node[box, below=of web] (cli) {CLI / Streamlit};
\node[box, right=of cli] (service) {Application service\\Parser, rules, scorer};
\node[box, right=of service] (reporter) {Reporter\\CSV/JSON artifact};
\node[box, below=of service] (logs) {CSV/JSON log\\+ config};
\draw[flow] (web) -- node[above, font=\scriptsize]{HTTP} (api);
\draw[flow] (api) -- (db);
\draw[flow] (cli) -- (service);
\draw[flow] (logs) -- (service);
\draw[flow] (service) -- (reporter);
\draw[flow] (service.east) to[bend left=18] node[above, font=\scriptsize]{permission check} (db.south);
\end{tikzpicture}
}
\caption{Kiến trúc RBAC Guard và demo Nova Bank.}
\label{fig:architecture}
\end{figure}

\section{Các thành phần chính}

Các thành phần của hệ thống được tổ chức thành bốn lớp trách nhiệm. Lớp trình
bày gồm giao diện Next.js, tiếp nhận thao tác của người dùng nhưng không tự đưa
ra quyết định phân quyền. Lớp nghiệp vụ và kiểm soát truy cập nằm tại FastAPI,
nơi xác thực phiên, kiểm tra permission, áp dụng chính sách nghiệp vụ và tạo
bản ghi audit cho mỗi hành động cần theo dõi. Lớp lưu trữ sử dụng SQLite cùng
\texttt{RBACRepository} để duy trì nhất quán giữa user, role, permission, phiên
đăng nhập và trạng thái giao dịch.

Ranh giới trách nhiệm này bảo đảm dữ liệu do trình duyệt gửi lên chỉ được xem
là yêu cầu, không phải bằng chứng về quyền hạn. FastAPI tra cứu quyền từ
repository trước khi thực thi thao tác; frontend chỉ hiển thị kết quả do API
trả về. Vì vậy, việc ẩn một nút trên giao diện không thay thế cho kiểm soát ở
backend, và một yêu cầu gọi trực tiếp endpoint vẫn phải đi qua cùng cơ chế xác
thực và phân quyền.

Ở nhánh giám sát, audit adapter chuyển các sự kiện xác thực và phân quyền trong
Nova Bank sang hợp đồng Event chung. Parser chuẩn hóa dữ liệu, tập luật phát
hiện chỉ báo, bộ chấm điểm gán mức rủi ro và reporter ghi các artifact phục vụ
điều tra. Chuỗi xử lý này tách schema nghiệp vụ khỏi logic phát hiện, đồng thời
cho phép truy vết từ một quyết định tại API tới cảnh báo hoặc incident tương
ứng. Bảng~\ref{tab:components} tổng hợp trách nhiệm cụ thể của từng thành phần.

\begin{table}[htbp]
\centering
\small
\begin{tabularx}{\textwidth}{p{0.23\textwidth}X}
\toprule
Thành phần & Trách nhiệm \\
\midrule
Next.js & Đăng nhập, quản trị nhân viên, thao tác giao dịch, xem phê duyệt và hiển thị phản hồi quyền. \\
FastAPI & Kiểm tra Bearer token, áp dụng dependency phân quyền, thực thi chính sách phê duyệt và ghi audit log. \\
SQLite/\\ RBACRepository & Lưu user, role, permission, session, giao dịch, chính sách demo và lịch sử hoạt động. \\
CLI/Service & Khởi tạo DB, kiểm tra quyền, điều phối một lần phân tích và đánh giá alert theo nhãn kỳ vọng. \\
Audit adapter & Chuyển bản ghi xác thực/phân quyền của Nova Bank thành Event chuẩn hóa mà không gắn nhãn đánh giá. \\
Parser và rules & Chuẩn hóa CSV/JSON, cô lập dòng lỗi và phát hiện SQL Injection, password guessing, unauthorized access. \\
Reporter & Ghi nguyên tử alert, metadata, context finding và incident dưới dạng CSV/JSON. \\
\bottomrule
\end{tabularx}
\caption{Phân chia trách nhiệm giữa các thành phần.}
\label{tab:components}
\end{table}

\section{Mô hình dữ liệu RBAC}

Mô hình dữ liệu được chia thành lớp Core RBAC và lớp hỗ trợ vận hành Nova Bank.
Lớp Core RBAC trong Hình~\ref{fig:rbac-core-erd} gồm năm quan hệ
\texttt{users}, \texttt{roles}, \texttt{permissions},
\texttt{user\_roles} và \texttt{role\_permissions}. Hai bảng liên kết chuyển
hai quan hệ nhiều--nhiều user--role và role--permission thành các quan hệ
một--nhiều có khóa ngoại rõ ràng. Cấu trúc này tránh gán permission trực tiếp
cho từng user và giữ role làm đơn vị quản trị chính sách.

\begin{figure}[htbp]
\centering
\resizebox{\textwidth}{!}{\begin{tikzpicture}[
  entity/.style={draw, rounded corners, align=left, font=\scriptsize,
    inner sep=2.5mm, minimum height=27mm, text width=27mm, fill=blue!5},
  link/.style={-Latex, thick},
  node distance=7mm
]
\node[entity] (users) {\textbf{users}\\
  \underline{id}\\username [UQ]\\status\\display\_name\\password\_hash};
\node[entity, right=of users] (userroles) {\textbf{user\_roles}\\
  \underline{user\_id} [FK]\\\underline{role\_id} [FK]\\PK(user\_id, role\_id)};
\node[entity, right=of userroles] (roles) {\textbf{roles}\\
  \underline{id}\\name [UQ]};
\node[entity, right=of roles] (rolepermissions) {\textbf{role\_permissions}\\
  \underline{role\_id} [FK]\\\underline{permission\_id} [FK]\\PK(role\_id, permission\_id)};
\node[entity, right=of rolepermissions] (permissions) {\textbf{permissions}\\
  \underline{id}\\resource\\action\\UQ(resource, action)};
\draw[link] (users) -- node[above, font=\scriptsize]{1:N} (userroles);
\draw[link] (roles) -- node[above, font=\scriptsize]{1:N} (userroles);
\draw[link] (roles) -- node[above, font=\scriptsize]{1:N} (rolepermissions);
\draw[link] (permissions) -- node[above, font=\scriptsize]{1:N} (rolepermissions);
\end{tikzpicture}
}
\caption{ERD của lớp Core RBAC. UQ ký hiệu ràng buộc duy nhất.}
\label{fig:rbac-core-erd}
\end{figure}

\begin{table}[htbp]
\centering
\small
\begin{tabularx}{\textwidth}{p{0.23\textwidth}p{0.25\textwidth}X}
\toprule
Quan hệ & Khóa và ràng buộc chính & Vai trò trong mô hình \\
\midrule
\texttt{users} & PK \texttt{id}; UQ \texttt{username}; \texttt{status} thuộc \texttt{active}/\texttt{disabled} & Lưu định danh, trạng thái, tên hiển thị và password hash của người dùng. \\
\texttt{roles} & PK \texttt{id}; UQ \texttt{name} & Đại diện chức năng nghiệp vụ như Administrator, Teller, Controller và Auditor. \\
\texttt{permissions} & PK \texttt{id}; UQ(\texttt{resource}, \texttt{action}) & Biểu diễn một quyền dưới dạng cặp tài nguyên--hành động. \\
\texttt{user\_roles} & PK(\texttt{user\_id}, \texttt{role\_id}); hai FK & Gán role cho user và ngăn một phép gán bị lặp. \\
\texttt{role\_permissions} & PK(\texttt{role\_id}, \texttt{permission\_id}); hai FK & Gán permission cho role và ngăn một quyền bị cấp lặp. \\
\bottomrule
\end{tabularx}
\caption{Từ điển dữ liệu của lớp Core RBAC.}
\label{tab:rbac-core-schema}
\end{table}

Khi API cần kiểm tra quyền, repository lần lượt nối \texttt{users} với
\texttt{user\_roles}, \texttt{role\_}\allowbreak\texttt{permissions} và
\texttt{permissions},
đồng thời lọc theo \texttt{username}, trạng thái
\texttt{active} và cặp \texttt{resource:action}, rồi chỉ cần xác nhận có ít
nhất một hàng phù hợp. \texttt{roles} không xuất hiện trực tiếp trong truy vấn
quyết định vì hai bảng liên kết đã chia sẻ \texttt{role\_id}; bảng này vẫn cần
thiết để quản trị tên role và hiển thị ma trận quyền. Backend thực hiện truy vấn
này trước hành động được bảo vệ, do đó việc frontend ẩn nút chỉ cải thiện trải
nghiệm và không phải là cơ chế phân quyền.

Trong Nova Bank, \texttt{transactions:read\_own} giới hạn Teller ở giao dịch
do chính họ tạo, còn \texttt{transactions:read\_all} cho phép Controller,
Auditor hoặc Administrator đọc phạm vi rộng hơn. Permission
\texttt{transactions:approve} chỉ được seed cho Controller. Đây là chính sách
Maker--Checker cụ thể ở API, không phải mô hình Separation of Duty tổng quát;
schema hiện tại chưa có bảng constraint để biểu diễn xung đột nhiệm vụ.

\section{Tích hợp audit log với pipeline}

Lệnh \texttt{export-audit-events} nhận database Nova Bank và tạo file CSV có
cùng schema với log đầu vào. Một bản ghi login được ánh xạ thành Event
\texttt{authentication}; các quyết định permission, policy và nghiệp vụ được
ánh xạ thành \texttt{authorization}. Outcome \texttt{denied} được giữ nguyên
để \texttt{UnauthorizedAccessRule} có thể kết hợp dấu vết ứng dụng với ma trận
quyền hiện tại. Adapter không ghi \texttt{expected\_label}, vì dữ liệu tích hợp
không được dùng làm ground truth.

Thiết kế này tạo một seam rõ ràng thay vì để detector phụ thuộc trực tiếp vào
schema nội bộ của Nova Bank. Nếu sau này audit schema bổ sung client IP, request
ID hoặc service name, adapter có thể mở rộng phép ánh xạ mà không đổi rule.
Giới hạn hiện tại là source IP dùng sentinel \texttt{nova-bank-local}; do đó
không nên dùng audit export này để đánh giá các luật phụ thuộc IP cho tới khi
API thu thập địa chỉ nguồn đáng tin cậy.
